\section{A test that can fail within a reader's lifetime}
\label{sec:testing}

\subsection{A dated workforce claim with a stable accounting boundary}

The primary dated hypothesis is \textbf{80\% fewer people working in the DGF ecosystem by
20 September 2033 than in 2026}. Its proposed operational test uses a registered population of
organizations and their DGF suppliers, with a stable definition of DGF activities. Let $P_y$ be the
mean of twelve monthly headcounts for the year ending on 20 September of year $y$. For $P_{2026}>0$,
the threshold is
\begin{equation}
P_{2033}/P_{2026}\le 0.20. \label{eq:workforce2033}
\end{equation}
Thus the prediction is tested as a reduction of at least 80\%, not a residual workforce of 80\%.
The final observation window runs from 20 September 2032 to 20 September 2033; baseline records cover
the corresponding twelve months ending in 2026. This averaging rule measures a sustained reduction
rather than a single favorable census date. The protocol is proposed: no enterprise cohort has been
enrolled and no preregistration identifier is claimed. Missing historical baseline records prevent
confirmation of this dated test; they cannot be replaced silently by a later baseline.

Each monthly count includes every person who performed in-scope DGF work during that month, counted
once across the registered ecosystem, regardless of employer or contract type. Internal staff,
contractors, external reviewers, provider operators, automation supervisors, maintenance, and recovery
workers remain in scope. Shared provider staff are deduplicated across participating customers;
their assigned FTE and hours are allocated separately. A worker performing even part-time DGF work
remains in the headcount. Privacy-preserving linkage and auditable supplier coverage are prerequisites;
unobserved supplier labor is missing data, not zero people. Hours and assigned FTE are additional
outcomes, not substitutes for the people count.

The organizational cohort and activity definitions are fixed before evaluation, but new DGF work and
new roles within that boundary remain included. Outsourcing, renaming teams, moving work between
participating firms, or adding agent-support jobs cannot make workers disappear from the account.
Mergers and exits are tracked under prespecified continuity rules; lost coverage cannot be scored as
a workforce reduction. Genuine redeployment outside DGF is a population exit, not necessarily a loss
of employment. A cohort test supports inference to the wider DGF ecosystem only with justified
sampling, coverage, and weighting; a convenience sample cannot establish an ecosystem-wide percentage.

In parallel, record project volume, case mix, quality, service levels, effective authority, and changes
in demand or scope. Report raw headcount and workload-adjusted human hours separately. The headcount
threshold could be reached through lower demand or organizational closure; attributing it to AI requires
a comparative design that addresses those alternatives. Automation-based substitution also requires
preserved quality and service within prespecified margins and valid authority.

Complete execution closure and zero recurring human support remain stronger, separate research
outcomes without a fixed calendar deadline. One unresolved contract or support task can coexist with
the 80\% reduction. Each registered contract still specifies evidence, dispositions, quality, authority,
and its full human-work boundary. Appendix~\ref{app:contract} supplies an illustrative contract.

\subsection{Measurement before deployment and after delegation}

The first study can use historical logs across all nine populations. Apply a versioned shadow-routing
rule using only information that would have been available at routing time. Keep later effort
measurements hidden from the rule to avoid selecting expensive cases by construction. With $J$ the
assigned residual and $T$ baseline effort, estimate $q$, $\kappa$, and $\phi=\mathbb{E}[TJ]/\mathbb{E}[T]$,
with uncertainty and stratification by route. This identifies where human work concentrates without
deploying a single agent; it does not measure $\eta$, automation quality, or workforce replacement.

The next stage compares human, assisted, and no-human execution on adjudicated cases, including valid
refusals, missing evidence, and attempted manipulation. A no-human run counts as replacement only when
quality, authorized action, and downstream acceptance are met. Research evaluators establish quality
independently; continuing evaluator work required for operation counts in the support ledger. Ordinary
production time, exceptional investigation, after-hours repair, supplier labor, external quality
review, and fixed maintenance have disjoint activity codes. Delays are recorded separately from active
hours.

\subsection{Producer-by-format experiment}

The handoff experiment crosses upstream producer with artifact format (Table~\ref{tab:factorial}),
stratified by route. It assigns eligible cases before production and measures upstream and downstream
human effort, including rejected outputs. Downstream tools and substantive policies are common within
each route. Related integration work packages are clustered by integration program to account for
shared evidence and reviewers.

\begin{table}[!htb]
\centering\small
\caption{Four conditions separating production from presentation.}
\label{tab:factorial}
\begin{tabularx}{0.92\linewidth}{@{}P{3.6cm}LL@{}}
\toprule
 & Free-form artifact & Structured, evidence-bearing artifact\\
\midrule
Human production & Human narrative/dossier & Human-authored typed record\\
Agent-assisted production & Agent narrative/dossier & Agent-authored typed record\\
\bottomrule
\end{tabularx}
\end{table}

Two estimands are distinct. A \emph{format-only} experiment fixes adjudicated facts and randomizes
their presentation, measuring the effect of representation. A \emph{production} experiment assigns
producer and format before facts are extracted, retaining omissions, inconsistencies, and verification
failures as outcomes. Restricting analysis to correct agent outputs would select away part of the
treatment cost. For case $c$ in cluster $j$, estimate separate outcomes $v\in\{h,c\}$ (hours or cost):
\begin{equation}
Y^{(v)}_{cj}=\beta^{(v)}_{0}+\beta^{(v)}_{A}A_{cj}+\beta^{(v)}_{F}F_{cj}+\beta^{(v)}_{AF}A_{cj}F_{cj}
+\gamma^{(v)\top}x_{cj}+u_j+e_{cj}. \label{eq:factorial}
\end{equation}
Here $A$ denotes agent-assisted production and $F$ structured format. For total human-hours, H2a
predicts $\beta^{(h)}_{F}<0$; the format effect among agents is $\beta^{(h)}_{F}+\beta^{(h)}_{AF}$.
For upstream production-and-verification cost, H2b predicts $\beta^{(c)}_{A}+\beta^{(c)}_{AF}<0$,
including tools and allocated assurance. Downstream quality, labor, and rework meet prespecified
equivalence margins across structured conditions. All assigned cases remain in the analysis;
coefficients are outcome-specific. This prevents simple standardization from being attributed to agent
intelligence.

Shared teams, caches, and process learning can create spillovers. Cluster or time-block randomization
can address them, with carryover specified in the analysis. Human time is measured directly or checked
against sampled observation; entry-to-exit elapsed time includes waiting and is not an effort measure.
Process mining helps reconstruct visits and routes \citep{vanderaalst2016}. Failed, abandoned, and
censored cases remain in the reporting boundary.

\subsection{Prospective tests of the six hypotheses}

For H1, two independent assessors score readiness before agent outcomes are available. Prediction is
evaluated on held-out cases or units against simpler duration, completeness, and model-configuration
baselines. Assessor agreement and probability calibration are reported separately.

For H3, the shadow-routing study above identifies $\kappa$, $\phi$, and their uncertainty within each
route; substantive exceptions are separated from random audits or fixed authorization rules.
Post-deployment effort on that partition identifies $\eta$. The frozen partition distinguishes baseline
difficulty from additional recovery work caused by deployment. Variance and post-treatment duration
are secondary outcomes, not automatic implications of selection.

For H4, a training period estimates $q,\kappa,\eta,\mu,r,B$ and visit rates; predictions are frozen
for held-out cases or units. The uniform-effort comparator predicts
$\hat L_U=\hat\lambda\hat h_0[\hat q+(1-\hat q)\hat m]$ with $\hat m=\hat h_R/\hat h_0$, setting
conditional exception effort to the overall baseline and additional overhead to zero. Both models use
the same training period. Total hours include maintenance and shared assurance; a stratified case-mix
forecast supplies a second benchmark. Added parameters must improve held-out loss and
direction-of-change accuracy, not merely training fit. The discriminating test then changes the routing
rule, the case mix, or the agent configuration between training and evaluation: the component model
re-estimates $\phi$ on historical work under the new conditions and carries $\eta$ and $\mu$ over, while
the matched comparator carries over its conditional means $h_E$ and $h_R$. Ratios are estimated within
sufficiently homogeneous categories of cases, and the compared models receive the same information and
the same recalibration budget. A change in $\eta$ between a
historical period and a post-deployment period is attributed to the agent only when other differences,
in population, policy, or practice, are controlled; a frozen routing rule helps but does not remove them
all.

For H5, four records distinguish who produced analysis, executed tools, authorized the disposition,
and held answerability for it. Time-series adoption slopes or time-to-delegation contrasts are
evaluated where institutional arrangements allow variation. Immutable human-signature requirements are
recorded as fixed strata rather than counted as evidence for a discovered lag.

For H6, every residual human intervention is logged with a category and an origin: present at baseline,
displaced from another role or organization, or created by the automation (supervision, evaluation,
exception triage, agent maintenance). Each annual evaluation reports, per route, the number of open
categories, the hours they absorb, the categories closed, and the categories created. The taxonomy of categories is versioned: a category is closed only when a qualified capability executes
it under mandate on the declared case population, not when it is merely absent from the cases observed,
and renamed, merged, split, or displaced categories are traced so that the count cannot change without
a change in the work. Reopened categories are reported. In the terms of
Equation~\eqref{eq:renewal}, the registry estimates the closure rate $a$ and the creation rate $b$.
Complete substitution predicts that closures exceed creations until the registry is empty; durable
complementarity predicts a stable or renewing set. A registry can remain nonempty in 2033 while the workforce prediction succeeds.
It then supports partial substitution but does not establish complete closure. This is the test that separates the two
trajectories, and it is the one the law ultimately stands on.

A second, separate experiment compares no intervention, upstream only, downstream only, and both. It
estimates the interaction in Equation~\eqref{eq:interaction}. Positive factorial producer effects do
not themselves establish that deploying adjacent gates together is superadditive.

\subsection{Quality, adaptation, and the data boundary}

A minimum case record contains pseudonymous case, parent program or project, route, and visit
identifiers; population/activity codes; active hours; model/tool and policy versions; provenance;
routing reasons; adjudicated outcomes; returns; upkeep; and authorization references. Independent
assessors resolve contested cases under written criteria while retaining disagreement. Model
self-reported confidence is not treated as calibrated correctness.

Before evaluation, the analysis plan specifies the smallest relevant labor saving, noninferiority
margins for quality, forecast-error tolerance, clustering, sample-size assumptions, and stopping
conditions for unsafe approvals or disclosures. Rare harmful events require uncertainty bounds and
sufficient follow-up, not an inference of safety from zero observations. The strategic extension adds
independently verified compliance and a pre/post-policy analysis so that apparent progress cannot be
explained only by improved presentation of unchanged defects.

\subsection{Deadlines, plateaus, and the last human task}

Each annual evaluation reports the four terms of Equation~\eqref{eq:rho}, human minutes per visit,
separate support hours, adjudicated quality, and mandate status. A support cost divided among more
customers is reported separately from an actual reduction in support hours. At each route, the blocking
term is named: evidence pursuit, judgment, verification, per-case approval, maintenance, or legal
oversight.

A proposed companion trajectory test, to be preregistered, defines a plateau as a residual ratio above 5\%, with less
than 0.5 percentage points of reduction per year over three successive years, measured at comparable
output, quality, and evaluation resources. The thresholds are design choices, not natural constants.
Such a result rejects the continued-decline specification; it is warning evidence, not a logical
refutation of the 80\% workforce prediction. A plateau below 20\% of baseline headcount could satisfy
that prediction while contradicting continued movement toward zero.

At the deadline, a headcount ratio above 0.20 rejects the workforce prediction for the observed
cohort. With complete census coverage, apply the threshold directly; with sampling, report the
prespecified interval for the ratio, clustering by organization and accounting for shared suppliers.
An interval spanning 0.20 leaves the result inconclusive. Incomplete baseline or supplier coverage also
prevents confirmation. Report separately whether quality, service, and authority conditions hold and
whether the design supports attribution to AI. The twelve-month window captures periodic review and
maintenance; zero observed interventions would still not prove zero risk on arbitrary future inputs.

\subsection{What would defeat the law?}

An established indispensable human intervention defeats technical universality for its declared
domain. A legally required personal approval prevents complete adoption there while it remains
necessary. A provider-operated hidden review queue defeats an apparent no-human API, and recurring
human support defeats support closure. These findings need not reject an 80\% workforce reduction:
the claims have different thresholds.

The dated workforce prediction fails if the correctly measured ratio remains above 0.20 in 2033.
It cannot be rescued by dropping difficult sites, excluding new support jobs, weakening the people
count, or postponing the deadline. Meeting a lower workload requirement without a corresponding
headcount change does not confirm the prediction. Conversely, a headcount reduction alone does not
prove that AI caused it or that governance quality was preserved.

Publish actual people, assigned FTE, workload, qualified execution, effective mandates, and support
closure separately. The unrestricted statement ``one day, all gates'' has no ordinary finite-date
disproof; the 2033 workforce hypothesis supplies a refutable intermediate prediction without claiming
that its success would prove the universal endpoint.
